\subsection{Convolutional and Geometric Architectures}

The convolutional-network source is migrated here because convolution is an architectural prior for data with grid, group, or geometric structure.

\begin{definition} \label{def:sl-convolution} \hlt{Convolution}\\
For functions on the real line, convolution is
\[
(f*g)(t)=\int_{-\infty}^{\infty}f(\tau)g(t-\tau)\,d\tau.
\]
For sequences on $\Z$, the discrete convolution is
\[
(f*g)[n]=\sum_{m=-\infty}^{\infty}f[m]g[n-m].
\]
\end{definition}

\begin{definition} \label{def:sl-convolutional-neural-network} \hlt{Convolutional Neural Network}\\
A convolutional neural network uses convolution or cross-correlation in at least one layer.  Neural-network libraries often implement cross-correlation while referring to it as convolution.
\end{definition}

\begin{remark} \label{rem:sl-cnn-motivations} \hlt{Sparse Interactions, Sharing, and Equivariance}\\
Convolutional layers use small kernels, share parameters across spatial locations, and produce representations that are equivariant to translations.  Pooling then creates approximate invariance over local neighborhoods.
\end{remark}

\begin{definition} \label{def:sl-pooling-padding-stride} \hlt{Pooling, Padding, and Stride}\\
Pooling replaces nearby activations by summary statistics such as maxima or averages.  Padding controls boundary behavior, and stride controls the spacing between evaluated output locations.
\end{definition}

\begin{remark} \label{rem:sl-geometric-deep-learning} \hlt{Geometric Extensions}\\
Graph neural networks, equivariant networks, and manifold-based architectures generalize the convolutional idea: encode the symmetries and locality of the domain directly into the model class.
\end{remark}
